\chapter{Introducción}
\label{cap:introduccion}
\setcounter{page}{1}

Este capítulo presenta la idea general del trabajo, la motivación detrás del tema y los conceptos principales en los que se basa. Se habla de la electroencefalografía como fuente de información para detectar emociones, del reconocimiento automático de emociones y de la importancia de usar modelos que se puedan interpretar. Por último, se explica cómo está organizada la memoria.


\section{Contexto y motivación}

La electroencefalografía destaca por su capacidad para extraer datos de un órgano tan delicado como el cerebro sin necesidad de realizar una intervención invasiva. Esto la convierte en una herramienta muy valiosa en la neurología, ya sea para diagnosticar tumores y lesiones cerebrales, evaluar el deterioro cognitivo o incluso detectar emociones \cite{1}.

Cuando era pequeño me hicieron una prueba del sueño en la que me colocaron electrodos con una pasta adhesiva bastante difícil de quitar y me monitorizaron durante toda la noche. Me llamó mucho la atención: ¿cómo podía un médico mirar todos esos “garabatos” llamados señales y entender qué pasaba en mi cabeza?

Con el tiempo, esa curiosidad se juntó con mi interés por la inteligencia artificial y su uso en la medicina. Al ver cómo otras personas creaban modelos capaces de detectar tumores, me pregunté si también sería posible hacer algo parecido con electroencefalogramas. La idea no es sustituir a los médicos, sino crear una herramienta que pueda ayudarles a interpretar mejor grandes cantidades de datos y aportar información útil de forma automática y fiable.

Por ello, este trabajo se centra en el reconocimiento de emociones mediante señales EEG, dando importancia no solo a los resultados obtenidos, sino también a la interpretación de los modelos. En el ámbito biomédico, no basta con que un modelo clasifique bien: también es necesario entender qué información está utilizando y si sus decisiones pueden explicarse de forma clara.


\section{Electroencefalografía como fuente de información afectiva}

La electroencefalografía (EEG) permite medir la actividad eléctrica del cerebro colocando electrodos sobre el cuero cabelludo. Esto no muestra directamente lo que ocurre en cada neurona, ofrece información sobre la actividad cerebral general, la cual puede cambiar según el estado mental o emocional de una persona. Por ello, estas señales se han convertido en una fuente de información relevante para el estudio de emociones.

A diferencia de otras señales fisiológicas, el EEG proporciona información directamente relacionada con la actividad cerebral, lo que resulta especialmente interesante en el ámbito de la computación afectiva. Las emociones no siempre se manifiestan de forma clara a través de expresiones faciales, voz o comportamiento externo, por lo que el análisis de señales cerebrales puede aportar una perspectiva complementaria y menos dependiente de respuestas observables.

Sin embargo, trabajar con EEG presenta importantes dificultades. Las señales son ruidosas, tienen baja amplitud y pueden verse afectadas por movimientos, parpadeos, actividad muscular o diferencias entre sujetos. Además, la relación entre una emoción concreta y un patrón cerebral no suele ser directa ni evidente. Por este motivo, es necesario aplicar técnicas de preprocesamiento, extracción de características y aprendizaje automático que permitan identificar patrones relevantes dentro de la señal.

En este trabajo, el EEG se considera una fuente de información afectiva porque permite explorar la relación entre la actividad cerebral y los estados emocionales. Su análisis mediante modelos automáticos ofrece la posibilidad de clasificar emociones a partir de datos fisiológicos, manteniendo el interés en comprender qué componentes de la señal resultan más importantes para dicha clasificación.


\section{Reconocimiento automático de emociones}
\label{sec:reconocimiento_automatico_emociones}

% Presentar el problema de clasificar emociones automáticamente y situar las tres
% condiciones consideradas en el trabajo: alegría, tristeza y estado neutro.

\section{IA de caja blanca e interpretabilidad en señales biomédicas}
\label{sec:ia_caja_blanca_interpretabilidad}

% Justificar por qué en señales biomédicas no basta con obtener una buena métrica,
% sino que también interesa comprender qué variables, canales o bandas utiliza el modelo.

\section{Estructura de la memoria}
\label{sec:estructura_memoria}

% Resumir brevemente qué contiene cada capítulo de la memoria.
